% ─────────────────────────────────────────────────────────────────────────────
% Template «manuale» — preambolo.
% Non caricare babel e non toccare tocdepth: se ne occupa opzioni.tex.
% ─────────────────────────────────────────────────────────────────────────────

\usepackage[T1]{fontenc}
\usepackage{avant}
\usepackage{helvet}
\usepackage{courier}
\renewcommand{\familydefault}{\sfdefault}
% avant e helvet si contendono \sfdefault: vince l'ultimo caricato, quindi il
% corpo resta in Helvetica e Avant Garde si richiama a mano nei titoli.
\newcommand{\fonttitolo}{\fontfamily{pag}\selectfont}

\usepackage{geometry}
\geometry{left=24mm,right=24mm,top=20mm,bottom=22mm,headsep=7mm,footskip=12mm}

\usepackage{xcolor}
\definecolor{accento}{HTML}{C8571B}
\definecolor{grafite}{HTML}{201C19}
\definecolor{fumo}{HTML}{6E645C}
\definecolor{nebbia}{HTML}{F6F0E9}

\usepackage{graphicx}
\usepackage{booktabs}
\usepackage{array}
\usepackage{enumitem}
% Le procedure si leggono per passi: gli elenchi numerati sono l'ossatura del
% documento, non un ornamento, e stanno larghi.
\setlist[enumerate]{itemsep=6pt,parsep=2pt,topsep=6pt,leftmargin=*}
\setlist[itemize]{itemsep=3pt,parsep=0pt,topsep=5pt,leftmargin=*}

\usepackage{listings}
\lstset{
  basicstyle=\ttfamily\footnotesize,
  backgroundcolor=\color{nebbia},
  frame=single, rulecolor=\color{accento!45},
  framesep=8pt, xleftmargin=4pt, framexleftmargin=4pt,
  keywordstyle=\color{accento}\bfseries,
  commentstyle=\color{fumo}\itshape,
  breaklines=true, showstringspaces=false
}

\usepackage{titlesec}
\titleformat*{\subsection}{\fonttitolo\normalsize\bfseries\color{accento}}
\titleformat*{\subsubsection}{\fonttitolo\normalsize\bfseries\color{grafite}}
\titlespacing*{\section}{0pt}{16pt}{7pt}
% La barra piena sul titolo di sezione è il segnale di «nuovo passo della
% procedura»: si deve riconoscere sfogliando, senza leggere.
% titlesec passa il titolo come argomento all'ultimo comando del blocco: si
% scrive \barratitolo senza #1, o TeX protesta per un parametro fuori posto.
\newcommand{\barratitolo}[1]{%
  \colorbox{accento}{\parbox{\dimexpr\textwidth-2\fboxsep}{%
    \vspace{2pt}\thesection\quad#1\vspace{2pt}}}}
\titleformat{\section}[block]{\fonttitolo\large\bfseries\color{white}}{}{0pt}{\barratitolo}

\usepackage{fancyhdr}
\pagestyle{fancy}
\fancyhf{}
\renewcommand{\headrulewidth}{0pt}
\fancyhead[L]{\footnotesize\color{fumo}\titolo}
\fancyhead[R]{\footnotesize\color{fumo}\versione}
\fancyfoot[L]{\footnotesize\color{fumo}\piede}
\fancyfoot[R]{\footnotesize\bfseries\color{accento}\thepage}

\linespread{1.10}
\setlength{\parskip}{5pt}
\setlength{\parindent}{0pt}

\usepackage[colorlinks=true,linkcolor=accento,urlcolor=accento,citecolor=accento]{hyperref}
\hypersetup{pdftitle={\titolo},pdfauthor={\autore}}
